\section{Marco teórico}
\label{sec:marco-teorico}

El marco teórico tiene dos partes. La primera revisa los trabajos que asignan elementos a posiciones para reducir la competencia entre vecinos y el marco de coloración de grafos que los unifica, base del modelo espacial. La segunda reúne las herramientas de probabilidad con que se modela la compra bajo incertidumbre: las distribuciones de conteo, el método de Monte Carlo y las cadenas de Markov.

\section{Trabajo relacionado}
\label{sec:trabajo-relacionado}

La distribución de especies pertenece a una familia de problemas que asignan elementos a posiciones de modo que se reduzca la competencia, la interferencia o la incompatibilidad entre vecinos. Esta sección revisa cinco trabajos de esa familia (\cref{subsec:competencia-literatura}) y el marco de coloración de grafos que los unifica (\cref{subsec:coloracion}).

\subsection{Reducción de competencia}
\label{subsec:competencia-literatura}

El trabajo más cercano al reto es el de \textcite{elizalde2025}, que asignan diez especies a las posiciones de siembra de un polígono siguiendo los lineamientos de CONAFOR (625 plantas por hectárea). El terreno es un grafo cuyos vértices son posiciones y cuyos arcos indican vecindad; una variable binaria indica qué especie ocupa cada posición, y un coeficiente de competencia por par de especies se contabiliza solo cuando ambas quedan en posiciones vecinas. Resuelven un modelo entero con CBC y un algoritmo genético de cromosoma entero. Con un límite de tres horas, el modelo exacto obtuvo soluciones factibles sin garantía de calidad entre 169 y 324 nodos y ninguna con 361, 400 y 625 nodos, mientras que el genético resolvió todos los tamaños (349.44~s en el mayor) con calidad igual o muy cercana a la del exacto donde hubo comparación.

\textcite{eisenblatter2002} estudian la asignación de frecuencias en redes celulares, donde la interferencia entre transmisores cumple el papel de la competencia entre plantas: cada transmisor es un vértice, los canales son colores y una penalización por par de transmisores mide la interferencia. Formulan el problema de mínima interferencia como un programa entero, muestran que es fuertemente NP-difícil y reportan que las relajaciones lineales dan cotas débiles, mientras que heurísticas como DSATUR seguidas de búsqueda local mejoran mucho los planes reales. \textcite{akplogan2011} asignan cultivos a unidades de terreno a lo largo de varios años como un problema de satisfacción de restricciones ponderadas; con bloques interdependientes, la instancia de 120 unidades no se cerró en 48 horas. \textcite{liu2026} planean cultivos con una matriz de adyacencia $W$ entre unidades vecinas y una matriz de interacción $M$ entre cultivos, positiva para pares complementarios y negativa para pares competitivos, e incorporan incertidumbre en rendimientos y precios con escenarios Monte Carlo y optimización robusta. Por último, \textcite{kobal2025} muestran en bosques mixtos que los índices de competencia que consideran la distancia a los vecinos explican mejor el crecimiento que los que la ignoran (el $R^{2}$ ajustado pasa de 0.325 sin competencia a 0.647 con el índice de Hegyi).

\begin{table*}[t]
  \centering
  \caption{Comparación de los trabajos sobre reducción de competencia.}
  \label{tab:comparacion-competencia}
  \small
  \begin{tabularx}{\textwidth}{@{}>{\RaggedRight\arraybackslash}p{2.7cm}>{\RaggedRight\arraybackslash}p{3.1cm}*{2}{>{\RaggedRight\arraybackslash}X}@{}}
    \toprule
    Referencia & Dominio & Criterio de competencia o incompatibilidad & Método \\
    \midrule
    \parencite{elizalde2025} & Reforestación & Coeficiente por par de especies en posiciones vecinas & Programación entera y algoritmo genético \\
    \parencite{eisenblatter2002} & Telecomunicaciones & Interferencia por par de transmisores (0 a 1) & Programación entera, heurísticas y cotas \\
    \parencite{akplogan2011} & Agricultura & Efecto del cultivo previo, compatibilidad con el suelo y agrupación espacial & Restricciones ponderadas con ramificación y acotamiento \\
    \parencite{liu2026} & Agricultura bajo incertidumbre & Matriz de interacción entre cultivos vecinos & Optimización robusta con escenarios Monte Carlo \\
    \parencite{kobal2025} & Manejo forestal & Índices según tamaño y distancia de los vecinos & Regresión con optimización de umbrales \\
    \bottomrule
  \end{tabularx}
\end{table*}

La \cref{tab:comparacion-competencia} resume los cinco trabajos. De ellos se desprenden tres lecciones para el reto. Primero, la forma dominante de medir la competencia es un coeficiente por par de elementos vecinos en un grafo \parencite{elizalde2025,eisenblatter2002,liu2026}, y la distancia importa \parencite{kobal2025}, lo que respalda limitar la competencia a los vecinos inmediatos de la retícula. Segundo, los modelos exactos pierden escalabilidad rápidamente \parencite{elizalde2025,akplogan2011}, lo que justifica las metaheurísticas \parencite{elizalde2025,eisenblatter2002}; un modelo exacto sigue siendo útil en instancias pequeñas para medir qué tan lejos queda la heurística del óptimo. Tercero, solo \textcite{liu2026} incorporan incertidumbre, y ninguno estima la cantidad de plantas a adquirir cuando la supervivencia es incierta, que es la primera parte del reto.

\subsection{Coloración de grafos con penalización}
\label{subsec:coloracion}

Sea $G=(V,E)$ un grafo no dirigido. Una coloración propia con $k$ colores es una función $c:V\to\{1,\dots,k\}$ tal que $c(u)\neq c(v)$ para toda arista $uv\in E$; el menor $k$ para el que existe es el número cromático $\chi(G)$, y encontrarlo es NP-difícil \parencite{eisenblatter2002}. Como la coloración propia no basta para modelar requisitos prácticos, existen generalizaciones que, en lugar de prohibir los conflictos, minimizan una penalización total sobre las aristas \parencite{eisenblatter2002}, y sus variantes se han extendido a problemas de programación de tareas y horarios \parencite{dewerra2015}.

La distribución de especies se representa con un grafo cuyos vértices son las posiciones de siembra, cuyas aristas unen posiciones vecinas, cuyos colores son las especies y en el que cada arista pesa la competencia entre las especies asignadas a sus extremos. Esta representación coincide con la de \textcite{elizalde2025} y con la matriz de adyacencia de \textcite{liu2026}. En el tresbolillo, cada posición interior tiene seis vecinas equidistantes, de modo que el grafo es una retícula triangular y todas las aristas representan la misma distancia.

Una coloración propia no sirve directamente, por dos razones: no respeta la cantidad de plantas de cada especie, y trata la competencia como todo o nada, cuando en realidad es parcial y distinta para cada par de especies (incluido el par de una especie consigo misma). El modelo de mínima interferencia de \textcite{eisenblatter2002} y el de \textcite{elizalde2025} comparten, en cambio, la misma estructura: variables binarias de asignación, una restricción que asigna exactamente un color a cada vértice y una penalización por arista que se activa según los colores de sus extremos, cuya suma se minimiza. Esta coloración con penalización y conteo fijo por color es el marco que adopta la distribución espacial de este trabajo.

\subsection{Herramientas para modelar la incertidumbre}
\label{sec:marco-estocastico}

La compra se decide antes de saber cuántas plantas nuevas sobrevivirán. Esta sección reúne las distribuciones con que se modelan los conteos del problema y el método de Monte Carlo con el que se evalúan las decisiones bajo incertidumbre.

\subsubsection{Distribuciones de conteo}
\label{subsec:distribuciones}

\paragraph{Binomial} Si se plantan $x$ plantas de una especie y cada una sobrevive, de forma independiente, con probabilidad $q$, el número de sobrevivientes es $S(x)\sim\operatorname{Bin}(x,q)$, con
\begin{equation}
  \mathbb{P}(S(x)=k)=\binom{x}{k}q^{k}(1-q)^{x-k},\quad k=0,\dots,x
  \label{eq:binomial}
\end{equation}
y media $xq$ \autocite{:berger2002}. Comprar $\lceil r/q\rceil$ plantas para obtener $r$ sobrevivientes acierta solo en promedio: con esa compra, la meta se incumple con una probabilidad cercana a la mitad.

\paragraph{Binomial negativa} En una secuencia de ensayos de Bernoulli independientes con probabilidad de éxito $q$, el número de ensayos $X_r$ hasta obtener el $r$-ésimo éxito tiene distribución binomial negativa \autocite{:berger2002}:
\begin{equation}
  \mathbb{P}(X_r=x)=\binom{x-1}{r-1}q^{r}(1-q)^{x-r},\quad x=r,r+1,\dots
  \label{eq:binneg}
\end{equation}
En el problema, $X_r$ es el número de plantas que hay que plantar, una tras otra, hasta que $r$ sobrevivan. Las dos distribuciones se relacionan por la identidad
\begin{equation}
  \mathbb{P}\bigl(S(x)\geq r\bigr)=\mathbb{P}\bigl(X_r\leq x\bigr),
  \label{eq:dualidad}
\end{equation}
porque al menos $r$ de las primeras $x$ plantas sobreviven si y solo si la $r$-ésima sobreviviente aparece a más tardar en la planta $x$. En consecuencia, la menor compra que obtiene $r$ sobrevivientes con probabilidad al menos $1-\alpha$ es el cuantil $1-\alpha$ de $X_r$. Esta equivalencia convierte una restricción de probabilidad sobre la supervivencia en el cálculo de un cuantil.

\paragraph{Poisson} El número de ocurrencias de un evento en una región, cuando estas ocurren al azar con intensidad constante $\lambda$ por unidad de área, sigue una distribución de Poisson, $\mathbb{P}(E=k)=e^{-\lambda A}(\lambda A)^{k}/k!$ para una región de área $A$ \autocite{:berger2002}. Es el modelo natural para el número de plantas existentes por polígono si estas se dispersaron al azar. Su media y su varianza son iguales; cuando los conteos muestran una varianza mayor que la media (sobredispersión), la alternativa habitual es una binomial negativa vista como mezcla de Poisson.

\subsubsection{Método de Monte Carlo}
\label{subsec:montecarlo}

El método de Monte Carlo estima una cantidad $\theta=\mathbb{E}[g(\mathbf{X})]$ promediando evaluaciones del modelo sobre $K$ realizaciones independientes de la entrada aleatoria \autocite{:robert2004}:
\begin{equation}
  \hat{\theta}_K=\frac{1}{K}\sum_{k=1}^{K}g(\mathbf{X}_k)\xrightarrow{\text{c.s.}}\theta,
  \label{eq:montecarlo}
\end{equation}
donde la convergencia la garantiza la ley fuerte de los grandes números. Cada realización, o escenario, es un valor completo de las variables aleatorias del modelo; se obtiene transformando números pseudoaleatorios uniformes en la distribución deseada \autocite{:ross2013} y sobre él el modelo se evalúa de forma determinista. Una probabilidad es el caso $g=\mathbf{1}\{\cdot\}$, de modo que $\mathbb{P}(S(x)\geq r)$ se estima con la fracción de escenarios en que la compra $x$ alcanza la meta, y un cuantil se estima con el estadístico de orden correspondiente de los escenarios.

Si $\sigma^{2}=\operatorname{Var}[g(\mathbf{X})]<\infty$, el teorema central del límite da el intervalo aproximado de nivel $1-\alpha$
\begin{equation}
  \hat{\theta}_K\pm z_{\alpha/2}\frac{S_K}{\sqrt{K}},
  \label{eq:intervalo}
\end{equation}
con $S_K$ la desviación estándar muestral; para que el semiancho no exceda un error $\varepsilon$ se requieren $K\geq(z_{\alpha/2}S_K/\varepsilon)^{2}$ escenarios \autocite{:ross2013}. El error decrece como $K^{-1/2}$ sin importar la dimensión del problema: reducirlo a la mitad cuesta cuatro veces más escenarios. El método admite modelos no lineales arbitrarios y entrega la distribución completa de la salida, no solo su media, pero sus resultados no son mejores que las distribuciones supuestas para las entradas \autocite{:robert2004}.

\subsubsection{Cadenas de Markov}
\label{subsec:markov}

Un proceso $\{X_t\}_{t\geq0}$ con espacio de estados finito $\mathcal{X}$ es una cadena de Markov homogénea si la probabilidad de pasar al estado $j$ depende solo del estado actual \autocite{:ross2014}:
\begin{equation}
  \mathbb{P}(X_{t+1}=j\mid X_t=i,X_{t-1},\dots,X_0)=\mathbb{P}(X_{t+1}=j\mid X_t=i)=P_{ij}.
  \label{eq:markov}
\end{equation}
Las probabilidades de transición forman una matriz estocástica $\mathbf{P}$, con $P_{ij}\geq0$ y $\sum_jP_{ij}=1$, y si $\boldsymbol{\pi}_0$ es la distribución inicial, la distribución tras $t$ pasos es $\boldsymbol{\pi}_0\mathbf{P}^{t}$ \autocite{:norris1997}. Una distribución $\boldsymbol{\pi}$ con $\boldsymbol{\pi}\mathbf{P}=\boldsymbol{\pi}$ es estacionaria: describe la fracción de tiempo que la cadena pasa en cada estado.

En el problema hay dos cadenas naturales. La primera describe el orden de plantación: si cada planta es un estado, su especie, $P_{ij}$ es la probabilidad de que la siguiente planta del recorrido sea de la especie $j$ cuando la actual es de la $i$. Un arreglo que reduce la competencia tiene $P_{ii}=0$, porque nunca planta dos individuos de la misma especie seguidos, y su distribución estacionaria es la composición del arreglo. La segunda describe la vida de cada posición de la retícula, con estados vacía, plántula y establecida: una plántula muere o se establece, y una posición vacía solo cambia si se replanta. Con la distribución inicial dada por el conteo, esta cadena estimaría cuántas posiciones quedarán vacías en cada campaña de reposición. Este trabajo usa la primera para describir las claves (\cref{subsec:simulacion}) y se limita al horizonte de un año en la segunda, con la supervivencia binomial de la \cref{eq:binomial}.

